\section*{Bab \ref{ch:g10:freshwater-resources} --- Sumber Daya Air Tawar}

\begin{solution}{exo:g10:freshwater-resources:1}
\omterm{def:g10:freshwater-resources:stock}{Stok} adalah jumlah air yang tersimpan dalam sebuah reservoir pada suatu waktu,
dalam \unit{km^3} (atau kilogram). \omterm{def:g10:freshwater-resources:stock}{Fluks} adalah jumlah yang berpindah dari satu
reservoir ke reservoir lain per satuan waktu, dalam \unit{km^3} per tahun.
\omterm{def:g10:freshwater-resources:residence-time}{Waktu tinggal} adalah waktu rata-rata yang dihabiskan air dalam sebuah
reservoir, dalam hari atau tahun; dalam keadaan tunak \omterm{def:g10:freshwater-resources:residence-time}{waktu tinggal} sama
dengan \omterm{def:g10:freshwater-resources:stock}{stok} dibagi \omterm{def:g10:freshwater-resources:stock}{fluks}.
\end{solution}

\begin{solution}{exo:g10:freshwater-resources:2}
Sekitar \qty{2.5}{\percent} dari seluruh air adalah air tawar. Dari air tawar
itu, \qty{68.7}{\percent} membeku dalam tudung es, gletser, dan salju abadi,
dan \qty{30.1}{\percent} berada di bawah tanah. Sungai menampung paling
sedikit, \num{2120}~\unit{km^3}.
\end{solution}

\begin{solution}{exo:g10:freshwater-resources:3}
\omterm{def:g10:freshwater-resources:renewable}{Air tawar terbarukan} adalah air tawar yang diperbarui setiap tahun oleh
presipitasi dan mengalir melalui sebuah wilayah sebagai air larian dan air
tanah: volume per tahun, yaitu \omterm{def:g10:freshwater-resources:stock}{fluks}. \omterm{def:g10:freshwater-resources:fossil-groundwater}{Air tanah fosil} adalah air yang
memasuki akuifer sejak lama dan tidak diperbarui oleh imbuhan masa kini:
air itu adalah simpanan yang hanya dapat dikuras, yaitu \omterm{def:g10:freshwater-resources:stock}{stok}.
\end{solution}

\begin{solution}{exo:g10:freshwater-resources:4}
$91\,000/2120 \approx 43$ kali lebih banyak.
\end{solution}

\begin{solution}{exo:g10:freshwater-resources:5}
Atmosfer memperbarui airnya paling cepat (sekitar \num{9.5} hari), samudra
paling lambat di antara reservoir yang dihitung (sekitar \num{3200} tahun); air
tanah dalam lebih tua lagi. $3200 \times 365.25/9.5 \approx \num{120000}$
kali lebih lama.
\end{solution}

\begin{solution}{exo:g10:freshwater-resources:6}
Masuk: $413\,000 + 73\,000 = \num{486000}$~\unit{km^3} setahun, jadi
$T = 12\,700/486\,000 \approx \num{0.0261}$ tahun $\approx \num{9.5}$ hari.
Keluar: $373\,000 + 113\,000 = \num{486000}$~\unit{km^3}, sama besarnya: udara
berada dalam keadaan tunak, dan \omterm{def:g10:freshwater-resources:residence-time}{waktu tinggalnya} sama.
\end{solution}

\begin{solution}{exo:g10:freshwater-resources:7}
$1\,335\,040/413 \approx \num{3230}$ tahun. Dalam keadaan tunak samudra
menerima sebanyak yang hilang darinya: $373 + 40 = 413$ ribu~\unit{km^3}
setahun, sehingga \omterm{def:g10:freshwater-resources:residence-time}{waktu tinggal} yang dihitung dengan masukan sama besarnya.
\end{solution}

\begin{solution}{exo:g10:freshwater-resources:8}
Daratan menerima \num{113000}~\unit{km^3} setahun; \num{73000} kembali ke
udara melalui penguapan dan transpirasi dan \num{40000} mengalir ke laut:
$73\,000 + 40\,000 = \num{113000}$. Bagian yang kembali ke udara adalah
$73/113 \approx \qty{65}{\percent}$.
\end{solution}

\begin{solution}{exo:g10:freshwater-resources:9}
$273.2 \times 10^6 \times 1233.5 \approx \qty{3.37e11}{m^3}$, yaitu
\qty{337}{km^3}: hampir sebanyak seluruh air yang masih ditampung Laut Aral
pada tahun 1987.
\end{solution}

\begin{solution}{exo:g10:freshwater-resources:10}
$15.8 \times 0.3048 \approx \qty{4.82}{m}$. Selama 65 tahun dari 1950 sampai
2015, $482/65 \approx \qty{7.4}{cm}$ setahun.
\end{solution}

\begin{solution}{exo:g10:freshwater-resources:11}
$2/0.01 = 200$ tahun. \omterm{def:g10:freshwater-resources:stock}{Stok} kini turun sebesar $0.05 - 0.01 =
\qty{0.04}{km^3}$ setahun, sehingga \omterm{def:g10:freshwater-resources:stock}{stok} itu bertahan $2/0.04 = 50$ tahun.
\end{solution}

\begin{solution}{exo:g10:freshwater-resources:12}
1960: $1093/71.1 \approx 15$ tahun; 2010: $98.1/11.4 \approx 8.6$ tahun.
Penguapan mengambil lapisan air yang kira-kira sama tebalnya setiap tahun dari
setiap kilometer persegi danau; ketika danau menyusut, luasnya menyusut,
begitu pula volume total yang diuapkan.
\end{solution}

\begin{solution}{exo:g10:freshwater-resources:13}
(1) Kehilangan bersihnya $1.2 - 0.2 = \qty{1.0}{m}$ setahun di atas
$100 \times 10^6$~\unit{m^2}: $10^8$~\unit{m^3}, \qty{0.1}{km^3} setahun harus
mengalir masuk. (2) Aliran masuk turun menjadi \qty{0.05}{km^3}: danau itu
kehilangan $\qty{0.05}{km^3} = 5 \times 10^7$~\unit{m^3}, lapisan setebal
$5 \times 10^7/10^8 = \qty{0.5}{m}$. (3) Ketika muka air turun, tepian
danau muncul dari air dan luas danau menyusut: permukaan lebih kecil,
penguapan lebih sedikit, sampai penguapan dan aliran masuk kembali seimbang.
\end{solution}

\begin{solution}{exo:g10:freshwater-resources:14}
Akuifer itu menampung $100 \times 0.20 = \qty{20}{m}$ air, yaitu
\qty{20000}{mm}: dengan \qty{1}{mm} setahun, \num{20000} tahun. Air dalam
akuifer Sahara berumur \num{200000} sampai satu juta tahun, sepuluh sampai lima
puluh kali lebih tua: air itu diimbuh selama iklim masa lalu yang lebih basah.
\end{solution}

\begin{solution}{exo:g10:freshwater-resources:15}
$337/453\,000 \approx \qty{0.00074}{km}$, sekitar \qty{0.74}{m} air. Muka air
turun \qty{4.8}{m}, sekitar enam setengah kali lebih besar. Air hanya mengisi
pori-pori batuan (dan tidak seluruhnya terkuras keluar): bila muka air tanah
turun satu meter, air yang dilepaskan adalah lapisan yang jauh lebih tipis
daripada satu meter, di sini sekitar $0.74/4.8 \approx \num{0.15}$~m.
\end{solution}

\begin{solution}{pb:g10:freshwater-resources:1}
\textbf{1.} $1093/68\,000 \approx \qty{0.016}{km}$, sekitar \qty{16}{m}.

\textbf{2.} $1093 \times 10^9 \approx \qty{1.09e12}{m^3}$.

\textbf{3.} \qty{10}{g/L} sama dengan \qty{10}{kg/m^3}: $1.093 \times 10^{12} \times 10
\approx \qty{1.1e13}{kg}$, sekitar \num{11} miliar ton.

\textbf{4.} $1093/71.1 \approx 15$ tahun.

\textbf{5.} Tidak ada sungai yang keluar darinya: airnya hanya dapat pergi
sebagai uap (atau dengan merembes ke dalam tanah), sedangkan garamnya tetap
tinggal.

\textbf{6.} $71.1/68\,000 \approx \qty{0.00105}{km}$, sekitar \qty{1.05}{m}.

\textbf{7.} $9.4/68\,000 \approx \qty{0.00014}{km}$, sekitar \qty{0.14}{m}.

\textbf{8.} $71.1 - 9.4 = \qty{61.7}{km^3}$ setahun.

\textbf{9.} $79 + 37 = \qty{116}{km^3}$: $61.7/116 \approx \qty{53}{\percent}$,
sekitar separuhnya.

\textbf{10.} Masukan $1.0 + 6.2 = \qty{7.2}{km^3}$; keluaran \qty{36.8}{km^3};
perubahan bersih $7.2 - 36.8 = \qty{-29.6}{km^3}$.

\textbf{11.} $(502.7 - 327.2)/5 = \qty{35.1}{km^3}$ setahun: orde besarnya sama
dengan \qty{29.6}{km^3} pada neraca tahun 1987 (tahun-tahun tunggal berbeda:
penguapan \qty{47.9}{km^3} pada tahun 1984).

\textbf{12.} $36.8/0.00105 \approx \qty{35000}{km^2}$.

\textbf{13.} $35\,000/68\,000 \approx \qty{52}{\percent}$, sekitar separuhnya.

\textbf{14.} Danau itu telah menyusut: separuh permukaan menguapkan kira-kira
separuh air.

\textbf{15.} $345.6/1093 \approx \qty{32}{\percent}$, sekitar sepertiga.

\textbf{16.} $10 \times 1093/345.6 \approx \qty{31.6}{g/L}$.

\textbf{17.} Salinitas yang terukur lebih rendah: sebagian garam telah
meninggalkan air, mengkristal di dasar yang tersingkap, dan terbawa angin
sebagai debu bergaram.

\textbf{18.} Sekitar \qty{1.05}{m} setahun, yaitu tebal air yang diuapkan.

\textbf{19.} $7.2/35\,000 \approx \qty{0.0002}{km}$, sekitar \qty{0.20}{m}.

\textbf{20.} $1.05 - 0.20 \approx 0.85$ (\num{0.84} dengan nilai yang tidak
dibulatkan): \textbf{pada tahun 1987 muka air Laut Aral turun sekitar
\qty{0.8}{m}, mendekati satu meter dalam satu tahun saja}.
\end{solution}
